% Sebastian Trying to edit this file on 2026-03-08
% Guideline & Policy for Project Proposal
% Only one project proposal needs to be submitted per group. The idea behind this proposal is
% that I can give you feedback and suggestions to make you come up with a feasible plan for the
% project. You don’t have to do exactly what you say in your proposal and can even completely
% change your solving strategies afterwards if you want. But it’s important to have at least one
% reasonable plan/baseline to start from. The length should be 4 to 6 pages, not including
% appendices.


% Structure
% It should contain title, authors, abstract, introduction, related work, model/method, and
% experiments. In particular, all these sections should be complete (though they can be short)
% except for model/method and experiments.
\documentclass{article}

% Use final option for non-anonymous class submission
\usepackage[final]{neurips_2021}

% Recommended packages
\usepackage[utf8]{inputenc}
\usepackage[T1]{fontenc}
\usepackage{hyperref}
\usepackage{url}
\usepackage{booktabs}
\usepackage{amsfonts}
\usepackage{nicefrac}
\usepackage{microtype}
\usepackage{xcolor}
\usepackage{amsmath}
\usepackage{graphicx}

\title{Spring 2026 Analytics II Team Project Proposal Report}

% List the group and authors
\author{
  Group 1 \\[1ex]  % Adds Group 1 at the top
  Joshua \\
  Author 2 \\
  Author 3 \\
  Author 4 \\
  Author 5
}

% Set the date manually (March 8, 2026)
%\date{March 8, 2026}

%\author{
%  First Author \\
%  Department of XYZ \\
%  University Name \\
%  City, Country \\
%  \texttt{author1@university.edu}
%  \And
%  Second Author \\
%  Department of XYZ \\
%  University Name \\
%  City, Country \\
%  \texttt{author2@university.edu}
%}


\begin{document}

\maketitle

% Add professor's name below title page (separate from authors)
\vspace{2em}  % Adds vertical space between the authors and professor's name
\begin{center}
  \textbf{Professor:} K. Hua
\end{center}

% Abstract (15%)
% It should summarize the main idea of the project and its contributions. It should be
% understandable to anyone in the course.
\newpage
\begin{abstract}
Briefly summarize the problem, the proposed method, and the main findings.
Keep it concise (typically 150–250 words).
\end{abstract}



% Introduction (15%)
% It should clearly state the problem being addressed and why it is important.
\section{Introduction}
Introduce the research problem, motivation, and contributions of the paper.



% Related Work (15%)
% It should clearly distinguish what your proposed contribution from previous literature (e.g., a
% 1-2 sentence summary of at least 3 closely related papers).
\section{Related Work}
Discuss prior research and explain how your work differs from or extends it.



% Model/Method (25%)
% 1. It should contain the description of your research idea. You can use equations or sketches
% to help explain the idea.

% 2. It should contain a list of to-dos (having some time estimation could be helpful) that could
% in principle make you finish the project. Here’s an example:
% • Extended literature search and finishing related work section (3 hours)
% • Finding appropriate datasets and converting them to same format (2 hours)
% • Downloading xxx baseline and reproducing results of prior work (4 hours)
% • Implementing the proposed model (4 hours)
% • Running baselines on datasets and tweaking designs and or hyperparameters (4 hours)
% • Making figures and tables (2 hours)
% • Writing the project report (4 hours)

% You should almost always start from simple ideas and gradually add the complexity since
% a research idea typically has to FAIL many times until it works. It is highly likely that the
% project doesn’t go as planned. Therefore, with simple ideas, you could have the flexibility to
% tweak and a higher chance of obtaining some conclusions in the end.
\section{Model/Method}
Describe your model, algorithm, or methodology.

Example equation:
\begin{equation}
    y = f(x; \theta)
\end{equation}



% Experiments Design (20%)
% It should contain a list of to-do experiments (e.g., each experiment takes one subsection). One
% should clearly describe reasons of the experiments. One should also put placeholder figures and
% tables (with specifically designed axes, rows, columns, and so on). At last, one should write
% down the goal and expected outcome of specific experiments.

% Typically, in machine learning and related areas, one should have at least one main experiment
% on a convincingly large real-world dataset to show that your proposed model/method
% empirically works better than prior work. Moreover, one usually adds other experiments on
% real-world datasets to analyze how important a specific model/method design is, i.e., the so-
% called ablation study. One can also design special synthetic experiments to showcase the
% benefits of the proposed model/method under certain (possibly extreme) conditions.
\section{Experiments Design}
Describe the experimental setup, datasets, evaluation metrics, and results.
Include figures and tables as needed.

\bibliographystyle{plain}
\bibliography{references}

\end{document}
